\begin{lemma}[\textup{(5.4)}: testing both weak forms against $w=u-v$]
Let $n\ge1$ and $p>1$.  Let $\Omega\subseteq\R^n$ be open, let $R>0$ with
$B_R(x_0)\subseteq\Omega$, let $u$ be a weak solution of
$\operatorname{div}A(x,u,\nabla u)=B(x,u,\nabla u)$ in $\Omega$ with weak
gradient $g$ (\noderef{IsWeakSolutionOn}), and let $(v,V,g_v)$ be a frozen
comparison datum for $(u,U,g)$ on $B_R(x_0)$ at the frozen field
$A^0=A(x_0,U(x_0),\cdot)$ (\noderef{FrozenComparisonOn}).  Assume
\[
A(\cdot,u,g),\ A^0(g),\ A^0(g_v)\in L^{p/(p-1)}\bigl(B_R(x_0)\bigr),
\qquad
g-g_v\in L^p\bigl(B_R(x_0)\bigr),
\qquad
B(\cdot,u,g)\in L^1\bigl(B_R(x_0)\bigr),
\]
let $T\ge0$ and suppose $|u-v|\le T$ almost everywhere on $B_R(x_0)$.  Then
\[
\int_{B_R(x_0)}\bigl(A^0(g)-A^0(g_v),\,g-g_v\bigr)dx
\ \le\
\Bigl|\int_{B_R(x_0)}\bigl(A(x,u,g)-A^0(g),\,g-g_v\bigr)dx\Bigr|
+\Bigl|\int_{B_R(x_0)}B(x,u,g)\,(u-v)\,dx\Bigr| .
\]
\end{lemma}
\begin{proof}
Write $B_R=B_R(x_0)$ and $w=u-v$, so $\nabla w=g-g_v$ in the sense that
$g-g_v$ is the weak gradient of $u-v$.  Note $|B_R|<\infty$ (Preamble: the
volume of a ball).

\emph{Step 1: the test pair.}  By clause (3) of \noderef{FrozenComparisonOn},
the pair $(v-u,\,g_v-g)$ belongs to $W^{1,p}_0(B_R)$
(\noderef{MemW1pZeroOn}).  Also $|v-u|=|u-v|\le T$ almost everywhere on $B_R$,
by hypothesis, and $T\ge0$.

\emph{Step 2: the equation on $\Omega$, tested against the bounded pair.}
Apply \noderef{WeakSolutionBoundedZeroTest} with the present $n$ and $p$, the
open set $\Omega$, the centre $x_0$, the radius $R$, the field $A$, the
right-hand side $B$, the pair $(u,g)$, the test pair
$(\varphi,g_\varphi)=(v-u,\,g_v-g)$ and the bound $T$.  Its hypotheses are:
$1\le n$ and $1<p$; $\Omega$ open; $0<R$; $B_R\subseteq\Omega$; the weak-solution
property \noderef{IsWeakSolutionOn}, a hypothesis of the present lemma;
$A(\cdot,u,g)\in L^{p/(p-1)}(B_R)$ and $B(\cdot,u,g)\in L^1(B_R)$, both
hypotheses; the membership $(v-u,g_v-g)\in W^{1,p}_0(B_R)$ from Step 1; $T\ge0$;
and the essential bound $|v-u|\le T$ from Step 1.  It gives
\[
\int_{B_R}\Bigl(\bigl(A(x,u,g),\,g_v-g\bigr)+B(x,u,g)\,(v-u)\Bigr)dx=0 .
\tag{1}
\]

\emph{Step 3: the frozen equation on $B_R$, tested against the same pair.}
Apply \noderef{WeakSolutionZeroTest} with the present $n$ and $p$, the open set
$B_R$, the field $(x,z,h)\mapsto A^0(h)$, the pair $(v,g_v)$ in the roles of its
$(u,g)$, and $(v-u,\,g_v-g)$ in the roles of its $(v,g_v)$.  Its hypotheses are:
$1\le n$ and $1<p$; $B_R$ open; the weak-solution clause
\[
\mathrm{IsWeakSolutionOn}\;n\;p\;B_R\;\bigl((x,z,h)\mapsto A^0(h)\bigr)\;
\bigl((x,z,h)\mapsto0\bigr)\;v\;g_v ,
\]
which is clause (1) of \noderef{FrozenComparisonOn} verbatim --- note that the
right-hand side there is identically $0$, which is exactly what that node
requires and is why the frozen half needs no bounded test function;
$A^0(g_v)\in L^{p/(p-1)}(B_R)$, a hypothesis of the present lemma; and
$(v-u,g_v-g)\in W^{1,p}_0(B_R)$ from Step 1.  It gives
\[
\int_{B_R}\bigl(A^0(g_v),\,g_v-g\bigr)dx=0 .
\tag{2}
\]

\emph{Step 4: subtracting, and why the integrals may be split.}  Each of the
four functions
\[
\bigl(A(x,u,g),g_v-g\bigr),\quad
\bigl(A^0(g),g_v-g\bigr),\quad
\bigl(A^0(g_v),g_v-g\bigr),\quad
B(x,u,g)\,(v-u)
\]
is integrable on $B_R$.  For the first three this is H\"older's inequality on
$B_R$ at the conjugate exponents $\frac p{p-1}$ and $p$ (Preamble): each pairs
a function in $L^{p/(p-1)}(B_R)$, by hypothesis, with
$g_v-g=-(g-g_v)\in L^p(B_R)$.  For the fourth,
$|B(x,u,g)(v-u)|\le T\,|B(x,u,g)|$ almost everywhere on $B_R$, and the
right-hand side is in $L^1(B_R)$ by hypothesis.  Consequently every integral
written below exists and the integrals may be added and subtracted term by
term.

Subtract (2) from (1):
\[
\int_{B_R}\bigl(A(x,u,g)-A^0(g_v),\,g_v-g\bigr)dx
+\int_{B_R}B(x,u,g)\,(v-u)\,dx=0 ,
\]
and replace $g_v-g$ by $-(g-g_v)$ and $v-u$ by $-(u-v)$ throughout, which
multiplies the whole identity by $-1$:
\[
\int_{B_R}\bigl(A(x,u,g)-A^0(g_v),\,g-g_v\bigr)dx
+\int_{B_R}B(x,u,g)\,(u-v)\,dx=0 .
\tag{3}
\]

\emph{Step 5: adding and subtracting $A^0(g)$.}  Split the first integral of (3)
along $A(x,u,g)-A^0(g_v)=\bigl(A(x,u,g)-A^0(g)\bigr)+\bigl(A^0(g)-A^0(g_v)\bigr)$,
legitimate by the integrability recorded in Step 4:
\[
\int_{B_R}\bigl(A^0(g)-A^0(g_v),\,g-g_v\bigr)dx
=-\int_{B_R}\bigl(A(x,u,g)-A^0(g),\,g-g_v\bigr)dx
-\int_{B_R}B(x,u,g)\,(u-v)\,dx .
\]
Bounding each term on the right by its absolute value, and using
$-a\le|a|$ twice, gives the assertion.  \emph{This is the thesis'
(5.4) of the thesis}, with the three displayed quantities its $I_1$, $I_2$ and
$I_3$; the identity at l.~1466--1469 that it is derived from is (3) above, and
it carries the $B$-integral, which is why Step 2 must test a \emph{bounded}
pair while Step 3 need not.
\end{proof}
